%% kapa: exam=1 round=32 year=2021 period=2 form=A subject=appraisal_law subject_ko=감정평가관계법규 seq=10 number=10
%% source: Q-Net 감정평가사 기출문제 (cst00309) artlSeq=5208397 file=extracted/1008/5208397/A형/2021년도 제32회 감정평가사 1차 2교시 A형.pdf
%% conversion: kordoc 4.24.51 -> tools/gen1.py (auto); review=proofread-sonnet; flags=none
\begin{question}{10}
국토의 계획 및 이용에 관한 법령상 조례로 따로 정할 수 있는 것에 해당하지 않는 것은? (단, 조례에 대한 위임은 고려하지 않음)
\begin{choices}
  \item 도시·군계획시설의 설치로 인하여 토지 소유권 행사에 제한을 받는 자에 대한 보상에 관한 사항
  \item 도시·군관리계획 입안시 주민의 의견 청취에 필요한 사항
  \item 대도시 시장이 지역여건상 필요하여 정하는 용도지구의 명칭 및 지정목적에 관한 사항
  \item 기반시설부담구역별 특별회계 설치에 필요한 사항
  \item 공동구의 점용료 또는 사용료 납부에 관한 사항
\end{choices}
\end{question}
\answer{1}
